% Options for packages loaded elsewhere
\PassOptionsToPackage{unicode}{hyperref}
\PassOptionsToPackage{hyphens}{url}
\documentclass[
  11pt,
  a4paper,
]{article}
\usepackage{xcolor}
\usepackage[margin=18mm]{geometry}
\usepackage{amsmath,amssymb}
\setcounter{secnumdepth}{-\maxdimen} % remove section numbering
\usepackage{iftex}
\ifPDFTeX
  \usepackage[T1]{fontenc}
  \usepackage[utf8]{inputenc}
  \usepackage{textcomp} % provide euro and other symbols
\else % if luatex or xetex
  \usepackage{unicode-math} % this also loads fontspec
  \defaultfontfeatures{Scale=MatchLowercase}
  \defaultfontfeatures[\rmfamily]{Ligatures=TeX,Scale=1}
\fi
\ifPDFTeX\else
  % xetex/luatex font selection
    \setmainfont[]{Noto Serif}
    \setsansfont[]{Noto Sans}
    \setmonofont[]{Noto Sans Mono}
  \setmathfont[]{Noto Sans Math}
\fi
% Use upquote if available, for straight quotes in verbatim environments
\IfFileExists{upquote.sty}{\usepackage{upquote}}{}
\IfFileExists{microtype.sty}{% use microtype if available
  \usepackage[]{microtype}
  \UseMicrotypeSet[protrusion]{basicmath} % disable protrusion for tt fonts
}{}
\makeatletter
\@ifundefined{KOMAClassName}{% if non-KOMA class
  \IfFileExists{parskip.sty}{%
    \usepackage{parskip}
  }{% else
    \setlength{\parindent}{0pt}
    \setlength{\parskip}{6pt plus 2pt minus 1pt}}
}{% if KOMA class
  \KOMAoptions{parskip=half}}
\makeatother
\ifLuaTeX
\usepackage[bidi=basic]{babel}
\else
\usepackage[bidi=default]{babel}
\fi
\babelprovide[main,import]{russian}
\ifPDFTeX
\else
\babelfont{rm}[]{Noto Serif}
\fi
% get rid of language-specific shorthands (see #6817):
\let\LanguageShortHands\languageshorthands
\def\languageshorthands#1{}
\setlength{\emergencystretch}{3em} % prevent overfull lines
\providecommand{\tightlist}{%
  \setlength{\itemsep}{0pt}\setlength{\parskip}{0pt}}
\usepackage{bookmark}
\IfFileExists{xurl.sty}{\usepackage{xurl}}{} % add URL line breaks if available
\urlstyle{same}
\hypersetup{
  pdftitle={Билеты по дискретным случайным процессам},
  pdflang={ru-RU},
  hidelinks,
  pdfcreator={LaTeX via pandoc}}

\title{Билеты по дискретным случайным процессам}
\author{}
\date{}

\begin{document}
\maketitle

\section{Билет 19. Случайные последовательности с независимыми
приращениями.
Блуждания}\label{ux431ux438ux43bux435ux442-19.-ux441ux43bux443ux447ux430ux439ux43dux44bux435-ux43fux43eux441ux43bux435ux434ux43eux432ux430ux442ux435ux43bux44cux43dux43eux441ux442ux438-ux441-ux43dux435ux437ux430ux432ux438ux441ux438ux43cux44bux43cux438-ux43fux440ux438ux440ux430ux449ux435ux43dux438ux44fux43cux438.-ux431ux43bux443ux436ux434ux430ux43dux438ux44f}

Источники: Ширяев 2; Сердобольская, лекции 2-3.

\subsection{1. Короткий
ответ}\label{ux43aux43eux440ux43eux442ux43aux438ux439-ux43eux442ux432ux435ux442}

Случайная последовательность \((S_n)_{n\ge0}\) имеет независимые
приращения, если для любых

\[
0\le n_0<n_1<\cdots<n_k
\]

случайные величины

\[
S_{n_1}-S_{n_0},\quad S_{n_2}-S_{n_1},\quad\ldots,\quad S_{n_k}-S_{n_{k-1}}
\]

независимы.

Главный дискретный пример - случайное блуждание:

\[
S_0=0,
\qquad
S_n=\xi_1+\cdots+\xi_n,
\]

где \(\xi_1,\xi_2,\ldots\) - независимые случайные величины. Если шаги
одинаково распределены, то приращения не только независимы, но и
стационарны:

\[
S_{n+k}-S_n\stackrel d=S_k.
\]

Для iid-шагов характеристическая функция суммы:

\[
\varphi_{S_n}(t)=\left(\varphi_\xi(t)\right)^n.
\]

Если \(\mathbb E\xi_1=a\) и \(\operatorname{Var}\xi_1=\sigma^2\), то

\[
\mathbb ES_n=na,
\qquad
\operatorname{Var}S_n=n\sigma^2.
\]

\subsection{2. Полный
ответ}\label{ux43fux43eux43bux43dux44bux439-ux43eux442ux432ux435ux442}

\subsubsection{19.1. Определение независимых
приращений}\label{ux43eux43fux440ux435ux434ux435ux43bux435ux43dux438ux435-ux43dux435ux437ux430ux432ux438ux441ux438ux43cux44bux445-ux43fux440ux438ux440ux430ux449ux435ux43dux438ux439}

Этот билет продолжает Билет 16 о случайных последовательностях и связан
с Билетом 09 о независимости. Здесь важный объект не просто
последовательность независимых величин, а процесс, у которого независимы
изменения на непересекающихся промежутках времени.

Пусть

\[
(S_n)_{n\ge0}
\]

\begin{itemize}
\tightlist
\item
  случайная последовательность со значениями в \(\mathbb R\),
  \(\mathbb Z\) или \(\mathbb R^d\). Приращением за промежуток от \(m\)
  до \(n\) называется
\end{itemize}

\[
S_n-S_m,
\qquad
0\le m<n.
\]

Последовательность имеет независимые приращения, если для любых моментов
времени

\[
0\le n_0<n_1<\cdots<n_k
\]

приращения по непересекающимся промежуткам

\[
S_{n_1}-S_{n_0},\quad
S_{n_2}-S_{n_1},\quad\ldots,\quad
S_{n_k}-S_{n_{k-1}}
\]

независимы.

Важно: независимые приращения не означают, что сами значения
\(S_0,S_1,S_2,\ldots\) независимы. Наоборот, значения обычно сильно
зависимы, потому что

\[
S_{n+1}=S_n+(S_{n+1}-S_n).
\]

Например, в случайном блуждании \(S_{10}\) содержит в себе все шаги,
которые входят в \(S_9\).

\subsubsection{19.2. Стационарные
приращения}\label{ux441ux442ux430ux446ux438ux43eux43dux430ux440ux43dux44bux435-ux43fux440ux438ux440ux430ux449ux435ux43dux438ux44f}

Если распределение приращения зависит только от длины промежутка, то
говорят о стационарных приращениях:

\[
S_{n+k}-S_n\stackrel d=S_k-S_0.
\]

Если обычно берут \(S_0=0\), то:

\[
S_{n+k}-S_n\stackrel d=S_k.
\]

\subsubsection{19.3. Конструкция процесса через суммы независимых
величин}\label{ux43aux43eux43dux441ux442ux440ux443ux43aux446ux438ux44f-ux43fux440ux43eux446ux435ux441ux441ux430-ux447ux435ux440ux435ux437-ux441ux443ux43cux43cux44b-ux43dux435ux437ux430ux432ux438ux441ux438ux43cux44bux445-ux432ux435ux43bux438ux447ux438ux43d}

Главная конструкция дискретного процесса с независимыми приращениями -
сумма независимых шагов. Пусть

\[
\xi_1,\xi_2,\ldots
\]

\begin{itemize}
\tightlist
\item
  независимые случайные величины. Определим
\end{itemize}

\[
S_0=0,
\qquad
S_n=\sum_{j=1}^n\xi_j.
\]

Тогда

\[
S_n-S_m=\sum_{j=m+1}^n\xi_j.
\]

Если промежутки времени не пересекаются, то соответствующие суммы
состоят из непересекающихся наборов независимых величин \(\xi_j\).
Поэтому такие приращения независимы.

Если \(\xi_j\) одинаково распределены, то приращение

\[
S_{n+k}-S_n=\xi_{n+1}+\cdots+\xi_{n+k}
\]

имеет то же распределение, что

\[
S_k=\xi_1+\cdots+\xi_k.
\]

Значит при iid-шагах процесс имеет независимые и стационарные
приращения.

\subsubsection{19.4. Характеристические функции и
моменты}\label{ux445ux430ux440ux430ux43aux442ux435ux440ux438ux441ux442ux438ux447ux435ux441ux43aux438ux435-ux444ux443ux43dux43aux446ux438ux438-ux438-ux43cux43eux43cux435ux43dux442ux44b}

Распределение суммы выражается через свертку. Если шаги iid с законом
\(\mu\), то

\[
S_n\sim \mu^{*n},
\]

где \(\mu^{*n}\) - \(n\)-кратная свертка меры \(\mu\). Через
характеристические функции это особенно просто. Пусть

\[
\varphi_\xi(t)=\mathbb E e^{it\xi_1}.
\]

Тогда из независимости:

\[
\varphi_{S_n}(t)
=
\mathbb E e^{it(\xi_1+\cdots+\xi_n)}
=
\prod_{j=1}^n\mathbb E e^{it\xi_j}
=
\left(\varphi_\xi(t)\right)^n.
\]

Для не одинаково распределенных независимых шагов:

\[
\varphi_{S_n}(t)=\prod_{j=1}^n\varphi_{\xi_j}(t).
\]

Если существуют первые два момента, то они складываются:

\[
\mathbb ES_n
=
\sum_{j=1}^n\mathbb E\xi_j,
\]

и при независимости

\[
\operatorname{Var}S_n
=
\sum_{j=1}^n\operatorname{Var}\xi_j.
\]

В iid-случае, если

\[
\mathbb E\xi_1=a,
\qquad
\operatorname{Var}\xi_1=\sigma^2,
\]

получаем

\[
\mathbb ES_n=na,
\qquad
\operatorname{Var}S_n=n\sigma^2.
\]

\subsubsection{19.5. Случайное
блуждание}\label{ux441ux43bux443ux447ux430ux439ux43dux43eux435-ux431ux43bux443ux436ux434ux430ux43dux438ux435}

Случайное блуждание - это процесс вида

\[
S_n=S_0+\xi_1+\cdots+\xi_n,
\]

обычно с \(S_0=0\) или заданным начальным состоянием. Если шаги
принимают значения в \(\mathbb Z\), то блуждание идет по целочисленной
решетке. Самый классический пример - простое блуждание на \(\mathbb Z\):

\[
P\{\xi_k=1\}=p,
\qquad
P\{\xi_k=-1\}=q=1-p.
\]

Если

\[
p=q=\frac12,
\]

блуждание называется простым симметричным.

Для простого блуждания

\[
S_n=\xi_1+\cdots+\xi_n.
\]

Если за \(n\) шагов было \(r\) шагов вправо и \(n-r\) шагов влево, то

\[
S_n=r-(n-r)=2r-n.
\]

Поэтому

\[
P\{S_n=x\}
=
\binom n{(n+x)/2}p^{(n+x)/2}q^{(n-x)/2},
\]

если \(n+x\) четно и \(|x|\le n\); иначе вероятность равна \(0\).

Для простого блуждания:

\[
\mathbb E\xi_1=p-q=2p-1,
\]

\[
\operatorname{Var}\xi_1=1-(p-q)^2=4pq.
\]

Следовательно,

\[
\mathbb ES_n=n(p-q),
\qquad
\operatorname{Var}S_n=4npq.
\]

В симметричном случае:

\[
\mathbb ES_n=0,
\qquad
\operatorname{Var}S_n=n.
\]

\subsubsection{19.6. Марковское свойство и
асимптотика}\label{ux43cux430ux440ux43aux43eux432ux441ux43aux43eux435-ux441ux432ux43eux439ux441ux442ux432ux43e-ux438-ux430ux441ux438ux43cux43fux442ux43eux442ux438ux43aux430}

Случайное блуждание является марковской цепью. Если текущее состояние
равно \(i\), то следующее состояние равно

\[
i+\xi_{n+1}.
\]

Будущий шаг независим от прошлого, поэтому условное распределение
будущего зависит только от текущего состояния. Для простого блуждания на
\(\mathbb Z\) переходные вероятности:

\[
p_{i,i+1}=p,
\qquad
p_{i,i-1}=q,
\]

а остальные переходы равны нулю. Это связывает билет с Билетом 22 и
Билетом 23.

Случайное блуждание обычно не является стационарным по значениям:
распределение \(S_n\) меняется с \(n\), дисперсия растет как \(n\). Но
при iid-шагах оно имеет стационарные приращения. Это типичная ловушка:
стационарные приращения не равны стационарности самого процесса.

Асимптотически блуждание связано с законом больших чисел и центральной
предельной теоремой. Если шаги iid и имеют конечное среднее \(a\), то

\[
\frac{S_n}{n}\xrightarrow{P} a
\]

по закону больших чисел. Если дисперсия \(\sigma^2>0\) конечна, то

\[
\frac{S_n-na}{\sigma\sqrt n}
\Rightarrow
N(0,1).
\]

Это связь с Билетом 21.

\subsection{3. Основные
определения}\label{ux43eux441ux43dux43eux432ux43dux44bux435-ux43eux43fux440ux435ux434ux435ux43bux435ux43dux438ux44f}

\subsubsection{Приращение}\label{ux43fux440ux438ux440ux430ux449ux435ux43dux438ux435}

Для последовательности \((S_n)\) приращение на промежутке \((m,n]\):

\[
S_n-S_m.
\]

\subsubsection{Независимые
приращения}\label{ux43dux435ux437ux430ux432ux438ux441ux438ux43cux44bux435-ux43fux440ux438ux440ux430ux449ux435ux43dux438ux44f}

Последовательность имеет независимые приращения, если для любых

\[
0\le n_0<n_1<\cdots<n_k
\]

величины

\[
S_{n_1}-S_{n_0},\ldots,S_{n_k}-S_{n_{k-1}}
\]

независимы.

\subsubsection{Стационарные
приращения}\label{ux441ux442ux430ux446ux438ux43eux43dux430ux440ux43dux44bux435-ux43fux440ux438ux440ux430ux449ux435ux43dux438ux44f-1}

Распределение приращения зависит только от длины промежутка:

\[
S_{n+k}-S_n\stackrel d=S_k-S_0.
\]

Если \(S_0=0\), то

\[
S_{n+k}-S_n\stackrel d=S_k.
\]

\subsubsection{Случайное
блуждание}\label{ux441ux43bux443ux447ux430ux439ux43dux43eux435-ux431ux43bux443ux436ux434ux430ux43dux438ux435-1}

Процесс, задаваемый суммами шагов:

\[
S_0=x,
\qquad
S_n=x+\sum_{j=1}^n\xi_j,
\]

где \(\xi_j\) - независимые шаги.

\subsubsection{\texorpdfstring{Простое блуждание на
\(\mathbb Z\)}{Простое блуждание на \textbackslash mathbb Z}}\label{ux43fux440ux43eux441ux442ux43eux435-ux431ux43bux443ux436ux434ux430ux43dux438ux435-ux43dux430-mathbb-z}

Блуждание с шагами

\[
\xi_j\in\{-1,1\}.
\]

Если

\[
P\{\xi_j=1\}=P\{\xi_j=-1\}=\frac12,
\]

оно называется простым симметричным блужданием.

\subsubsection{Блуждание с одинаково распределенными
шагами}\label{ux431ux43bux443ux436ux434ux430ux43dux438ux435-ux441-ux43eux434ux438ux43dux430ux43aux43eux432ux43e-ux440ux430ux441ux43fux440ux435ux434ux435ux43bux435ux43dux43dux44bux43cux438-ux448ux430ux433ux430ux43cux438}

Блуждание с одинаково распределенными шагами. В этом случае приращения
стационарны.

\subsection{4. Основные теоремы и
свойства}\label{ux43eux441ux43dux43eux432ux43dux44bux435-ux442ux435ux43eux440ux435ux43cux44b-ux438-ux441ux432ux43eux439ux441ux442ux432ux430}

\begin{itemize}
\tightlist
\item
  Если
\end{itemize}

\[
S_n=S_0+\sum_{j=1}^n\xi_j
\]

и \(\xi_j\) независимы, то \((S_n)\) имеет независимые приращения.

\begin{itemize}
\tightlist
\item
  Если шаги iid, то приращения независимы и стационарны:
\end{itemize}

\[
S_{n+k}-S_n\stackrel d=S_k-S_0.
\]

\begin{itemize}
\tightlist
\item
  Распределение суммы независимых шагов - свертка:
\end{itemize}

\[
P_{S_n}=\mu_1*\cdots*\mu_n.
\]

В iid-случае:

\[
P_{S_n}=\mu^{*n}.
\]

\begin{itemize}
\tightlist
\item
  Характеристические функции независимых шагов перемножаются:
\end{itemize}

\[
\varphi_{S_n}(t)=\prod_{j=1}^n\varphi_{\xi_j}(t).
\]

В iid-случае:

\[
\varphi_{S_n}(t)=\left(\varphi_\xi(t)\right)^n.
\]

\begin{itemize}
\tightlist
\item
  При существовании моментов:
\end{itemize}

\[
\mathbb ES_n=S_0+\sum_{j=1}^n\mathbb E\xi_j,
\]

\[
\operatorname{Var}S_n=\sum_{j=1}^n\operatorname{Var}\xi_j.
\]

\begin{itemize}
\tightlist
\item
  Случайное блуждание с независимыми шагами является марковской цепью.
\item
  Блуждание обычно не стационарно по значениям, но при iid-шагах имеет
  стационарные приращения.
\item
  Для iid-шагов применимы ЗБЧ и ЦПТ:
\end{itemize}

\[
\frac{S_n}{n}\xrightarrow{P}\mathbb E\xi_1,
\qquad
\frac{S_n-n\mathbb E\xi_1}{\sqrt{n\operatorname{Var}\xi_1}}
\Rightarrow N(0,1).
\]

\subsection{5. Примеры}\label{ux43fux440ux438ux43cux435ux440ux44b}

\subsubsection{Простое симметричное
блуждание}\label{ux43fux440ux43eux441ux442ux43eux435-ux441ux438ux43cux43cux435ux442ux440ux438ux447ux43dux43eux435-ux431ux43bux443ux436ux434ux430ux43dux438ux435}

Пусть

\[
P\{\xi_k=1\}=P\{\xi_k=-1\}=\frac12.
\]

Тогда

\[
S_n=\xi_1+\cdots+\xi_n
\]

и

\[
P\{S_n=x\}
=
\binom n{(n+x)/2}2^{-n},
\]

если \(n+x\) четно и \(|x|\le n\). Иначе вероятность равна \(0\).

Среднее и дисперсия:

\[
\mathbb ES_n=0,
\qquad
\operatorname{Var}S_n=n.
\]

\subsubsection{Несимметричное
блуждание}\label{ux43dux435ux441ux438ux43cux43cux435ux442ux440ux438ux447ux43dux43eux435-ux431ux43bux443ux436ux434ux430ux43dux438ux435}

Пусть

\[
P\{\xi_k=1\}=p,
\qquad
P\{\xi_k=-1\}=q=1-p.
\]

Тогда

\[
\mathbb ES_n=n(p-q),
\qquad
\operatorname{Var}S_n=4npq.
\]

Если \(p>q\), блуждание имеет положительный дрейф; если \(p<q\),
отрицательный.

\subsubsection{Сумма
Бернулли}\label{ux441ux443ux43cux43cux430-ux431ux435ux440ux43dux443ux43bux43bux438}

Пусть

\[
\xi_k\sim \operatorname{Bernoulli}(p),
\qquad
\xi_k\in\{0,1\}.
\]

Тогда

\[
S_n=\sum_{k=1}^n\xi_k
\]

считает число успехов за \(n\) испытаний:

\[
S_n\sim \operatorname{Bin}(n,p).
\]

Это процесс с независимыми стационарными приращениями, но он не является
блужданием на всей \(\mathbb Z\), а является неубывающим счетчиком.

\subsubsection{\texorpdfstring{Блуждание в
\(\mathbb R^d\)}{Блуждание в \textbackslash mathbb R\^{}d}}\label{ux431ux43bux443ux436ux434ux430ux43dux438ux435-ux432-mathbb-rd}

Пусть шаги \(\xi_k\) принимают значения в \(\mathbb R^d\). Тогда

\[
S_n=S_0+\xi_1+\cdots+\xi_n
\]

задает блуждание в многомерном пространстве. Например, на решетке
\(\mathbb Z^2\) можно на каждом шаге переходить к одному из соседей.

\subsubsection{Пуассоновский процесс как
аналог}\label{ux43fux443ux430ux441ux441ux43eux43dux43eux432ux441ux43aux438ux439-ux43fux440ux43eux446ux435ux441ux441-ux43aux430ux43a-ux430ux43dux430ux43bux43eux433}

Пуассоновский процесс \(N(t)\) в непрерывном времени имеет независимые и
стационарные приращения:

\[
N(t+s)-N(t)\sim \operatorname{Poisson}(\lambda s).
\]

Это не дискретная случайная последовательность, но полезный аналог: он
показывает, что идея независимых приращений используется и в непрерывном
времени.

\subsection{6. Схемы доказательств /
обоснований}\label{ux441ux445ux435ux43cux44b-ux434ux43eux43aux430ux437ux430ux442ux435ux43bux44cux441ux442ux432-ux43eux431ux43eux441ux43dux43eux432ux430ux43dux438ux439}

\subsubsection{Почему сумма независимых шагов имеет независимые
приращения}\label{ux43fux43eux447ux435ux43cux443-ux441ux443ux43cux43cux430-ux43dux435ux437ux430ux432ux438ux441ux438ux43cux44bux445-ux448ux430ux433ux43eux432-ux438ux43cux435ux435ux442-ux43dux435ux437ux430ux432ux438ux441ux438ux43cux44bux435-ux43fux440ux438ux440ux430ux449ux435ux43dux438ux44f}

Пусть

\[
S_n=\sum_{j=1}^n\xi_j.
\]

Для моментов

\[
0\le n_0<n_1<\cdots<n_k
\]

приращение

\[
S_{n_r}-S_{n_{r-1}}
=
\sum_{j=n_{r-1}+1}^{n_r}\xi_j.
\]

Разные приращения зависят от непересекающихся наборов независимых шагов.
Поэтому они независимы.

\subsubsection{Почему характеристическая функция суммы
факторизуется}\label{ux43fux43eux447ux435ux43cux443-ux445ux430ux440ux430ux43aux442ux435ux440ux438ux441ux442ux438ux447ux435ux441ux43aux430ux44f-ux444ux443ux43dux43aux446ux438ux44f-ux441ux443ux43cux43cux44b-ux444ux430ux43aux442ux43eux440ux438ux437ux443ux435ux442ux441ux44f}

Если \(\xi_j\) независимы, то

\[
\varphi_{S_n}(t)
=
\mathbb E\exp\left(it\sum_{j=1}^n\xi_j\right)
=
\mathbb E\prod_{j=1}^n e^{it\xi_j}.
\]

По независимости математическое ожидание произведения равно произведению
математических ожиданий:

\[
\varphi_{S_n}(t)
=
\prod_{j=1}^n\mathbb E e^{it\xi_j}
=
\prod_{j=1}^n\varphi_{\xi_j}(t).
\]

\subsubsection{Почему блуждание является марковской
цепью}\label{ux43fux43eux447ux435ux43cux443-ux431ux43bux443ux436ux434ux430ux43dux438ux435-ux44fux432ux43bux44fux435ux442ux441ux44f-ux43cux430ux440ux43aux43eux432ux441ux43aux43eux439-ux446ux435ux43fux44cux44e}

Для блуждания

\[
S_{n+1}=S_n+\xi_{n+1}.
\]

Шаг \(\xi_{n+1}\) независим от прошлого

\[
S_0,\ldots,S_n.
\]

Поэтому условное распределение \(S_{n+1}\) при известном прошлом зависит
только от \(S_n\). В счетном случае:

\[
P\{S_{n+1}=j\mid S_0=i_0,\ldots,S_n=i\}
=
P\{\xi_{n+1}=j-i\}.
\]

Это марковское свойство.

\subsubsection{Почему блуждание не стационарно по
значениям}\label{ux43fux43eux447ux435ux43cux443-ux431ux43bux443ux436ux434ux430ux43dux438ux435-ux43dux435-ux441ux442ux430ux446ux438ux43eux43dux430ux440ux43dux43e-ux43fux43e-ux437ux43dux430ux447ux435ux43dux438ux44fux43c}

Для симметричного блуждания

\[
\operatorname{Var}S_n=n.
\]

Значит распределение \(S_n\) зависит от \(n\). Например, \(S_0=0\) почти
наверное, а \(S_1\) принимает значения \(\pm1\). Поэтому значения
процесса не стационарны, хотя приращения стационарны.

\subsection{7. Что написать на
доске}\label{ux447ux442ux43e-ux43dux430ux43fux438ux441ux430ux442ux44c-ux43dux430-ux434ux43eux441ux43aux435}

\begin{enumerate}
\def\labelenumi{\arabic{enumi}.}
\tightlist
\item
  Приращение:
\end{enumerate}

\[
S_n-S_m.
\]

\begin{enumerate}
\def\labelenumi{\arabic{enumi}.}
\setcounter{enumi}{1}
\tightlist
\item
  Независимые приращения:
\end{enumerate}

\[
S_{n_1}-S_{n_0},\ldots,S_{n_k}-S_{n_{k-1}}
\quad\text{независимы}.
\]

\begin{enumerate}
\def\labelenumi{\arabic{enumi}.}
\setcounter{enumi}{2}
\tightlist
\item
  Блуждание:
\end{enumerate}

\[
S_0=0,
\qquad
S_n=\sum_{j=1}^n\xi_j.
\]

\begin{enumerate}
\def\labelenumi{\arabic{enumi}.}
\setcounter{enumi}{3}
\tightlist
\item
  Стационарные приращения при iid-шагах:
\end{enumerate}

\[
S_{n+k}-S_n\stackrel d=S_k.
\]

\begin{enumerate}
\def\labelenumi{\arabic{enumi}.}
\setcounter{enumi}{4}
\tightlist
\item
  Характеристическая функция:
\end{enumerate}

\[
\varphi_{S_n}(t)=\left(\varphi_\xi(t)\right)^n.
\]

\begin{enumerate}
\def\labelenumi{\arabic{enumi}.}
\setcounter{enumi}{5}
\tightlist
\item
  Моменты:
\end{enumerate}

\[
\mathbb ES_n=n\mathbb E\xi_1,
\qquad
\operatorname{Var}S_n=n\operatorname{Var}\xi_1.
\]

\begin{enumerate}
\def\labelenumi{\arabic{enumi}.}
\setcounter{enumi}{6}
\tightlist
\item
  Простое блуждание:
\end{enumerate}

\[
P\{\xi=1\}=p,
\qquad
P\{\xi=-1\}=q.
\]

\subsection{8. Типичные дополнительные
вопросы}\label{ux442ux438ux43fux438ux447ux43dux44bux435-ux434ux43eux43fux43eux43bux43dux438ux442ux435ux43bux44cux43dux44bux435-ux432ux43eux43fux440ux43eux441ux44b}

\textbf{Что является главным объектом билета?}

Случайная последовательность, у которой независимы приращения на
непересекающихся промежутках. Главный пример - случайное блуждание, то
есть сумма независимых шагов.

\textbf{В чем разница между независимыми значениями и независимыми
приращениями?}

Независимые значения означали бы независимость \(S_0,S_1,\ldots\).
Независимые приращения означают независимость разностей \(S_n-S_m\) на
непересекающихся промежутках. У блуждания значения обычно зависимы, но
приращения независимы.

\textbf{Что значит стационарность приращений?}

Это значит, что закон приращения зависит только от длины промежутка:

\[
S_{n+k}-S_n\stackrel d=S_k-S_0.
\]

При \(S_0=0\):

\[
S_{n+k}-S_n\stackrel d=S_k.
\]

\textbf{Почему iid-шаги дают независимые стационарные приращения?}

Независимость приращений следует из независимости непересекающихся
наборов шагов. Стационарность следует из одинакового распределения
шагов: сумма любых \(k\) последовательных шагов имеет один и тот же
закон.

\textbf{Как найти распределение \(S_n\)?}

В общем случае это свертка законов шагов:

\[
P_{S_n}=\mu_1*\cdots*\mu_n.
\]

В iid-случае:

\[
P_{S_n}=\mu^{*n}.
\]

Для простого блуждания получается биномиальная формула.

\textbf{Почему характеристическая функция равна степени?}

Потому что характеристическая функция суммы независимых величин равна
произведению характеристических функций:

\[
\varphi_{S_n}(t)
=
\prod_{j=1}^n\varphi_{\xi_j}(t).
\]

Если шаги одинаково распределены:

\[
\varphi_{S_n}(t)=\varphi_\xi(t)^n.
\]

\textbf{Почему блуждание является марковской цепью?}

Потому что

\[
S_{n+1}=S_n+\xi_{n+1},
\]

а новый шаг независим от прошлого. Поэтому при известном \(S_n\) прошлые
состояния не дают дополнительной информации о \(S_{n+1}\).

\textbf{Случайное блуждание стационарно?}

По значениям обычно нет: распределение \(S_n\) меняется с \(n\). Но при
iid-шагах у него стационарные приращения.

\textbf{Как связаны блуждания с ЗБЧ и ЦПТ?}

Если шаги iid, то

\[
\frac{S_n}{n}\xrightarrow{P}\mathbb E\xi_1
\]

по закону больших чисел, а при конечной ненулевой дисперсии

\[
\frac{S_n-n\mathbb E\xi_1}{\sqrt{n\operatorname{Var}\xi_1}}
\Rightarrow N(0,1).
\]

\textbf{Какая главная ошибка в этом билете?}

Путать независимость приращений с независимостью самих значений процесса
и путать стационарность приращений со стационарностью процесса.

\subsection{9. Типичные
ошибки}\label{ux442ux438ux43fux438ux447ux43dux44bux435-ux43eux448ux438ux431ux43aux438}

\begin{itemize}
\tightlist
\item
  Путать независимые значения и независимые приращения.
\item
  Считать, что случайное блуждание стационарно по значениям.
\item
  Забывать начальное состояние \(S_0\).
\item
  Забывать, что стационарность приращений требует одинакового
  распределения шагов.
\item
  Писать \(\varphi_{S_n}=\varphi_\xi^n\) без условия одинакового
  распределения шагов.
\item
  Складывать дисперсии без независимости шагов.
\item
  Путать блуждание на \(\mathbb Z\) и счетчик Бернулли: во втором случае
  шаги \(0/1\), а не \(\pm1\).
\item
  Считать марковское свойство полной независимостью состояний.
\end{itemize}

\subsection{10. Связи с другими
билетами}\label{ux441ux432ux44fux437ux438-ux441-ux434ux440ux443ux433ux438ux43cux438-ux431ux438ux43bux435ux442ux430ux43cux438}

\begin{itemize}
\tightlist
\item
  Билет 09: независимость и факторизация математических ожиданий
  используются для приращений и характеристических функций.
\item
  Билет 16: блуждание - случайная последовательность и случайный элемент
  пространства последовательностей.
\item
  Билет 17: у блуждания при iid-шагах стационарные приращения, но обычно
  не стационарные значения.
\item
  Билет 21: ЗБЧ и ЦПТ описывают асимптотику \(S_n\).
\item
  Билет 22: блуждание задает переходную матрицу \(p_{i,i+1}=p\),
  \(p_{i,i-1}=q\).
\item
  Билет 23: случайное блуждание - пример марковской цепи.
\end{itemize}

\subsection{11. Минимум на
удовлетворительно}\label{ux43cux438ux43dux438ux43cux443ux43c-ux43dux430-ux443ux434ux43eux432ux43bux435ux442ux432ux43eux440ux438ux442ux435ux43bux44cux43dux43e}

Нужно сказать:

\begin{enumerate}
\def\labelenumi{\arabic{enumi}.}
\tightlist
\item
  Приращение процесса:
\end{enumerate}

\[
S_n-S_m.
\]

\begin{enumerate}
\def\labelenumi{\arabic{enumi}.}
\setcounter{enumi}{1}
\item
  Независимые приращения: приращения на непересекающихся промежутках
  независимы.
\item
  Случайное блуждание:
\end{enumerate}

\[
S_0=0,
\qquad
S_n=\xi_1+\cdots+\xi_n,
\]

где \(\xi_k\) независимы.

\begin{enumerate}
\def\labelenumi{\arabic{enumi}.}
\setcounter{enumi}{3}
\tightlist
\item
  Если шаги iid:
\end{enumerate}

\[
\varphi_{S_n}(t)=\varphi_\xi(t)^n.
\]

\begin{enumerate}
\def\labelenumi{\arabic{enumi}.}
\setcounter{enumi}{4}
\tightlist
\item
  Главная ловушка: независимые приращения не означают независимые
  значения.
\end{enumerate}

\subsection{12. Минимум на
хорошо}\label{ux43cux438ux43dux438ux43cux443ux43c-ux43dux430-ux445ux43eux440ux43eux448ux43e}

Помимо минимума, нужно объяснить:

\begin{itemize}
\tightlist
\item
  что такое стационарные приращения;
\item
  почему сумма независимых шагов имеет независимые приращения;
\item
  как распределение суммы связано со сверткой;
\item
  почему характеристические функции перемножаются;
\item
  как считаются \(\mathbb ES_n\) и \(\operatorname{Var}S_n\);
\item
  почему случайное блуждание является марковской цепью;
\item
  разобрать простое симметричное блуждание.
\end{itemize}

\subsection{13. Что добавить для
отлично}\label{ux447ux442ux43e-ux434ux43eux431ux430ux432ux438ux442ux44c-ux434ux43bux44f-ux43eux442ux43bux438ux447ux43dux43e}

Для сильного ответа стоит добавить:

\begin{itemize}
\tightlist
\item
  формулу распределения простого блуждания:
\end{itemize}

\[
P\{S_n=x\}
=
\binom n{(n+x)/2}p^{(n+x)/2}q^{(n-x)/2};
\]

\begin{itemize}
\tightlist
\item
  условие четности \(n+x\);
\item
  различие между стационарностью процесса и стационарностью приращений;
\item
  переходные вероятности блуждания как марковской цепи;
\item
  связь с ЗБЧ и ЦПТ;
\item
  пример счетчика Бернулли и пуассоновского процесса как аналога.
\end{itemize}

\subsection{14. Устная версия ответа на 5
минут}\label{ux443ux441ux442ux43dux430ux44f-ux432ux435ux440ux441ux438ux44f-ux43eux442ux432ux435ux442ux430-ux43dux430-5-ux43cux438ux43dux443ux442}

\begin{enumerate}
\def\labelenumi{\arabic{enumi}.}
\tightlist
\item
  Приращение последовательности - это разность \(S_n-S_m\).
\item
  Процесс имеет независимые приращения, если приращения на
  непересекающихся промежутках независимы.
\item
  Главный пример - случайное блуждание:
\end{enumerate}

\[
S_0=0,
\qquad
S_n=\xi_1+\cdots+\xi_n.
\]

\begin{enumerate}
\def\labelenumi{\arabic{enumi}.}
\setcounter{enumi}{3}
\tightlist
\item
  Если шаги \(\xi_k\) независимы, то приращения независимы, потому что
  каждое приращение является суммой своего набора шагов.
\item
  Если шаги одинаково распределены, то приращения стационарны:
\end{enumerate}

\[
S_{n+k}-S_n\stackrel d=S_k.
\]

\begin{enumerate}
\def\labelenumi{\arabic{enumi}.}
\setcounter{enumi}{5}
\tightlist
\item
  Распределение \(S_n\) - свертка законов шагов.
\item
  Через характеристические функции:
\end{enumerate}

\[
\varphi_{S_n}(t)=\prod_{j=1}^n\varphi_{\xi_j}(t),
\]

а в iid-случае:

\[
\varphi_{S_n}(t)=\varphi_\xi(t)^n.
\]

\begin{enumerate}
\def\labelenumi{\arabic{enumi}.}
\setcounter{enumi}{7}
\tightlist
\item
  При существовании моментов:
\end{enumerate}

\[
\mathbb ES_n=n\mathbb E\xi_1,
\qquad
\operatorname{Var}S_n=n\operatorname{Var}\xi_1.
\]

\begin{enumerate}
\def\labelenumi{\arabic{enumi}.}
\setcounter{enumi}{8}
\tightlist
\item
  Для простого блуждания \(\xi_k=\pm1\), вероятности шагов \(p\) и
  \(q\).
\item
  Блуждание является марковской цепью, потому что следующий шаг
  независим от прошлого.
\item
  По значениям блуждание обычно не стационарно, но при iid-шагах имеет
  стационарные приращения.
\item
  Асимптотика описывается ЗБЧ и ЦПТ.
\end{enumerate}

\subsection{15. Если осталось 2 минуты перед
экзаменом}\label{ux435ux441ux43bux438-ux43eux441ux442ux430ux43bux43eux441ux44c-2-ux43cux438ux43dux443ux442ux44b-ux43fux435ux440ux435ux434-ux44dux43aux437ux430ux43cux435ux43dux43eux43c}

Запомнить:

\[
S_n-S_m
\]

\begin{itemize}
\tightlist
\item
  приращение.
\end{itemize}

Независимые приращения:

\[
S_{n_1}-S_{n_0},\ldots,S_{n_k}-S_{n_{k-1}}
\quad\text{независимы}.
\]

Блуждание:

\[
S_0=0,
\qquad
S_n=\sum_{j=1}^n\xi_j.
\]

При iid-шагах:

\[
S_{n+k}-S_n\stackrel d=S_k,
\qquad
\varphi_{S_n}(t)=\varphi_\xi(t)^n.
\]

Главные ловушки:

\begin{itemize}
\tightlist
\item
  независимые приращения не означают независимые значения;
\item
  стационарные приращения не означают стационарный процесс;
\item
  формула \(\varphi_{S_n}=\varphi_\xi^n\) требует iid-шагов.
\end{itemize}

\newpage

\end{document}
